\section{Introduction}
\label{sec:intro}

Large language model (LLM) agents are increasingly deployed not as
demonstrations but as production services, and in that transition they accrue
costs and risks that are invisible in a prototype. A long-running session
accumulates context: the transcript grows turn over turn, and stale tool
outputs---binary blobs, large intermediate documents---are carried forward
into every subsequent model call, inflating cost without adding information.
Credentials collected mid-session to authenticate an external action pass
through the model to reach the tool that needs them, and in doing so are
written into the transcript, the provider's request history, and every cache
the transcript touches. Making an agent durable, resumable, human-supervised,
and self-correcting requires further machinery still. Mainstream agent
frameworks leave each of these concerns to per-application code: the framework
supplies the model loop and the tool-calling plumbing, and the developer is
expected to prune context, isolate secrets, and wire up retry and
human-in-the-loop control by hand. Teams routinely omit this work, or
implement it inconsistently, because nothing in the framework requires it and
the failure is silent until it is expensive.

This paper studies Colmena, a declarative agent engine that takes a different
position on where these concerns belong. In Colmena an agent is not a program
but a declarative document describing a directed graph (DAG) of typed
nodes and edges; the engine, not the agent author, implements each node's
behavior. Because the agent is data the engine interprets rather than code the
engine merely hosts, the engine is in a position to enforce properties of every
agent it runs, uniformly and by construction, rather than trusting each
application to supply them.

Our thesis is that Colmena relocates production concerns---context lifecycle,
credential handling, human-in-the-loop, and retry---from per-application code
into invariants the engine enforces by construction, and that the declarative
DAG-as-data execution model is the enabler of that relocation. We do not claim
these concerns are unreachable in a conventional framework; we claim that
reaching them there is the developer's responsibility, discharged in
hand-written code that a team may omit or get wrong, whereas in Colmena they
hold by default.

We support this thesis with the following contributions.

\begin{itemize}
  \item We describe the Colmena execution model---the DAG-as-data agent, the
    typed node behaviors the engine supplies, and the resulting
    \emph{engine boundary} that separates engine-enforced behavior from
    application logic (\S\ref{sec:model}).
  \item We present a provider-authoritative evaluation methodology in which
    every framework's LLM calls are routed through a shared proxy, so that
    token counts and costs are read from the transcript that crossed the wire
    rather than from any figure a framework self-reports (\S\ref{sec:methodology}).
  \item We measure three relocated concerns against idiomatic-default
    competitor implementations. \emph{Context} becomes an engine-managed
    resource: over a fixed analyst session Colmena consumes
    $\approx$\num{39085} input tokens against \numrange{404095}{452359} for
    the competitor defaults, a gap a competitor closes only with hand-written
    per-tool scrubbing code (\S\ref{sec:context}). \emph{Credential isolation}
    holds by construction: plaintext secrets reach the model transcript in
    $0\%$ of Colmena runs against $100\%$ for every competitor default, while
    a hand-architected competitor arm also reaches $0\%$ (\S\ref{sec:credentials}).
    \emph{Production hardening}---durable human-in-the-loop, credential
    masking, automatic critic-retry, and secret isolation---is expressed as
    declarative configuration of the agent document rather than as
    hand-rolled application logic (\S\ref{sec:hardening}).
  \item We give an honest default-versus-code account of the comparison,
    reporting where it does not favor Colmena and stating plainly what the
    engine does, and does not, uniquely provide (\S\ref{sec:limitations}).
\end{itemize}

The remainder of the paper proceeds as follows. Section~\ref{sec:related-work}
situates Colmena among existing agent frameworks. Section~\ref{sec:model}
develops the execution model on which the results rest, and
Section~\ref{sec:methodology} the measurement basis they share.
Sections~\ref{sec:context}, \ref{sec:credentials}, and \ref{sec:hardening}
report the three concern results in turn. Section~\ref{sec:limitations}
records the limitations and negative results, and
Section~\ref{sec:conclusion} concludes.
